% !TEX root = ../main.tex
\chapter*{General Conclusion}
\addcontentsline{toc}{chapter}{General Conclusion}
\markboth{General Conclusion}{General Conclusion}

This project asked when a learned local radar detector actually provides an advantage over classical
CFAR, and what is required for that advantage to survive the move from controlled simulation to real
measurements. It answered that question in five stages, each of which changed the answer the
previous one would have given.

\section*{What the project built}

A simulator producing labelled FMCW range--Doppler maps with a calibrated cell-SNR definition; four
CFAR detectors on a measured window geometry; a \TinyParams{}-parameter convolutional detector whose
receptive field equals the CFAR window by construction; and one evaluation path shared by both arms
--- the same candidate extraction, the same one-to-one matching, the same operating-point
definition, the same statistics. Every comparison in this report is made through that single path,
which is what allows the differences to be attributed to the detectors rather than to the
measurement.

\section*{Answering the research questions}

\textbf{RQ1 --- what a learned local rule adds, in-domain.} On a fresh held-out synthetic test set,
at matched \emph{measured} false-alarm rate, the learned detector establishes gains in five challenge
conditions, ties in homogeneous clutter, and shows no significant difference on the
resolution-limited stress case after correction. It reaches $\Pd = 0.5$ at a cell SNR
\SI{\SynSnrGainDL}{\deci\bel} below CA-CFAR. This establishes potential under controlled conditions.
It establishes nothing about deployment.

\textbf{RQ2 --- where the advantage comes from.} Attribution by intervention rules out the cheap
explanation: giving CFAR the same integrated mainlobe energy makes it worse, not better, on every
condition. Guard-region access contributes a small tolerant-detection benefit and simultaneously
broadens the response, hurting exact centring. Most usefully, the frozen arm's localisation deficit
proved to be largely a consequence of its flat $3\times3$ supervision rather than an inherent
property of learned detection --- concentrating the label removes most of it at negligible detection
cost. Together these narrow the field to an adaptive learned decision rule without uniquely
identifying its mechanism, which the experiments were not designed to do.

\textbf{RQ3 --- whether it survives transfer.} It does not. Evaluated zero-shot on CARRADA, with no
retraining, fine-tuning or real-data recalibration, the synthetic-trained network and CA-CFAR are
\textbf{statistically indistinguishable} in exact dense-mask recall (\RoneTinyCNN{} against
\RoneCFARCA{}; sequence-balanced difference \RoneDelta{}, 95\,\% CI [\RoneCIlo{}, \RoneCIhi{}]),
while the network emits \RoneBurdenRatio{} times more unmatched detections per frame. The classical
detector does not reverse the result; the learned advantage simply disappears, and the learned arm
pays for parity with output density.

\textbf{RQ4 --- whether the training domain was the limiting factor.} Substantially, yes. The domain
audit measured large differences between the simulator and the real sensor in background smoothness,
spread, zero-Doppler ridge structure and range non-stationarity. Rebuilding the synthetic training
background from CARRADA \emph{training}-split statistics alone --- with no real frame entering
neural optimisation or threshold calibration --- raised exact test recall from \CiTestOrig{} to
\CiTestNew{} (\CiTestDelta{}, 95\,\% CI [\CiTestCIlo{}, \CiTestCIhi{}]), positive on every test
sequence, and closed the near-range deficit. But it did so at \CiTestBurdenRatio{} times the output
density, and burden-matched analysis showed CA-CFAR still ahead once output was equalised. Domain
alignment recovered recall; it did not by itself establish superiority.

\textbf{RQ5 --- what real supervision adds.} Fine-tuned on CARRADA training sequences, with the
operating point fixed on held-out training sequences before the evaluation split was opened, the
same network reached exact recall \FtValFt{} against CA-CFAR's \FtValCFAR{} on previously unscored
validation sequences --- a difference of \FtDelta{} (95\,\% CI [\FtCIlo{}, \FtCIhi{}]) at a slightly
\emph{lower} unmatched-detection burden (\FtValBurdenFt{} against \FtValBurdenCFAR{}). It remains
ahead under an evaluation-only burden sweep. Notably, fine-tuning established \emph{no} recall gain
over the domain-informed synthetic arm (\FtVsCiDelta{}, CI [\FtVsCiCIlo{}, \FtVsCiCIhi{}]); what it
changed was selectivity, cutting output roughly \FtBurdenDrop{}-fold at equal recall.

\section*{What the evidence supports}

\textbf{Recall alone would have been misleading at every stage.} Read on recall only, the five
stages say: the learned detector wins, then ties, then wins large, then wins. Read with output
burden beside it, they say something more useful and less flattering --- that the learned detector's
standing depends on how much output it is allowed to emit, and that the one stage where it is ahead
on both axes at once is the stage where it received real supervision.

\textbf{The fairness caveat that belongs with the final result.} The fine-tuned network is supervised
with the same dense segmentation support the primary metric scores against. That is legitimate
supervised adaptation --- it is what the experiment set out to test --- but it gives the learned
detector an opportunity the classical detector never receives. The defensible statement is that
\emph{under supervised target-domain adaptation, the learned detector outperformed CA-CFAR on this
evaluation}, not that it is inherently the better detector.

\textbf{The statistical caveat.} The final interval's lower bound is \FtCIlo{} on \FtNseq{}
sequences. It excludes zero, but narrowly; the fine-tuned arm is ahead on \FtSeqAhead{} of
\FtNseq{} sequences and marginally behind on the rest, and the aggregate leans on two sequences
where the classical detector does poorly. The result is real and it is not comfortable.

\section*{Limitations}

\begin{itemize}
\itemsep2pt
\item \textbf{One simulator, one real dataset, one architecture.} Nothing here supports a claim about
learned detectors in general.
\item \textbf{No false-alarm measurement on real data.} CARRADA annotates the objects that were
labelled, not every scatterer present, so the learned arm's denser output is an output-burden
measure and can never be priced against recall as a false-alarm rate.
\item \textbf{The domain intervention changed several properties jointly.} It supports a causal
statement about the combined intervention and about nothing narrower; no individual background
property has been isolated.
\item \textbf{Partial domain-gap reduction only.} The Doppler correlation structure was deliberately
left unmatched, because correcting it would have required changing the range--Doppler processing
window that also defines the target point spread function and the guard geometry.
\item \textbf{The attribution experiments run on validation data.} They explain the frozen benchmark;
they do not extend it, and no arm they identify as better was itself frozen and scored on held-out
data.
\item \textbf{The fine-tune was deliberately small and had not converged.} Every seed reached the
ten-epoch ceiling with its development loss still falling, so the final result is a lower bound on
that configuration rather than its best achievable behaviour.
\item \textbf{The post-processing protocol was revised during the work}, on synthetic development
data, and the real data re-scored under it. Appendix~\ref{app:prereg} records what changed.
\end{itemize}

\section*{Future work}

The results point to specific next experiments rather than general ambitions.

\textbf{Separate what the domain intervention bundled.} One ablation per property --- multi-look-like
smoothness, ridge modelling, range non-stationarity, correlation structure --- would convert a
combined causal statement into individual ones.

\textbf{Measure how much real data is actually needed.} The fine-tuning result used all fifteen
fitting sequences. A sample-efficiency curve over increasing fractions of real labelled data would
say whether the advantage requires a substantial annotation effort or a small one, which is the
question a practitioner actually faces.

\textbf{Prefer domain randomisation to calibration against one dataset.} The intervention here was
tuned toward a single sensor. Randomising background statistics over a plausible range would test
whether the recovered transfer generalises or is specific to CARRADA, and a second real dataset
would test the same thing directly.

\textbf{Make output burden part of the objective.} Every stage of this work found burden to matter as
much as recall, yet it was never optimised --- only measured. A burden-aware loss or calibration
criterion follows naturally.

\textbf{Adapt with weaker supervision.} Dense mask supervision is expensive. Whether weak,
self-supervised or unsupervised adaptation recovers any of the same benefit is open.

\textbf{Carry the improved supervision design through a complete benchmark.} The label ablation
showed concentrated supervision substantially improves localisation at negligible detection cost, but
that arm was never frozen, never scored on held-out data and never taken to real measurements.

Larger architectures are deliberately \emph{not} the first recommendation. The behaviour that
dominated every stage of this project --- domain sensitivity, operating-point dependence, the
recall-versus-burden trade --- was visible at \TinyParams{} parameters, and is better understood
there than hidden inside a larger model.

\section*{Final statement}

The results do not support replacing CFAR with deep learning. They support a conditional hierarchy.
CFAR is the stronger out-of-the-box detector: training-free, interpretable, cheap, and transferable
without adaptation --- it carried to a real sensor it had never seen with no change at all. A small
learned detector can exceed it, but its advantage is distribution-dependent. In this study the
controlled synthetic gains were lost under zero-shot transfer, substantially recovered by
target-informed synthetic modelling, and converted into a joint recall-and-efficiency advantage only
after supervised real-domain adaptation.

The practical question is therefore not whether deep learning beats CFAR, but whether enough domain
knowledge and representative training data are available to justify the learned detector's higher
performance ceiling. Where they are, it is worth the effort. Where they are not, CFAR is still the
right answer.
